\documentclass[11pt,a4paper]{article}
\usepackage[UTF8]{ctex}
\usepackage{geometry}
\geometry{left=2.5cm,right=2.5cm,top=2.5cm,bottom=2.5cm}
\usepackage{amsmath,amssymb,amsthm}
\usepackage{hyperref}
\usepackage{booktabs}
\usepackage{enumitem}
\usepackage{xltabular}
\hypersetup{colorlinks=true,linkcolor=blue,urlcolor=blue}

\newtheorem{definition}{定义}[section]
\newtheorem{proposition}{命题}[section]
\newtheorem{remark}{注记}[section]

\title{AI 化攻防批判（修正版）\\
——对称性博弈、战场转移与 AI 时代的够用就行}
\author{李龙胜}
\date{\today}

\begin{document}

\maketitle

\begin{center}
\textit{
当攻防双方都部署 AI 时，优势并不会自动归于防御方。\\
攻击方的 AI 按需生成，防御方却要全量分析海量告警——\\
但如果防御方将逃逸概率压到极低，攻击方的成本优势也会被逆转。\\
对抗样本军备竞赛将博弈从网络层精确转移到模型参数空间。\\
防御方面临虚警消耗、模型退化和分析师能力丧失的三重陷阱。\\
防御方在 AI 时代的建设性策略：\\
将模型选择、推理深度和对抗训练频率纳入贪心择优，\\
同时刻意保留人工研判以防止人的能力退化。\\
}
\end{center}

\vspace{5mm}

\begin{abstract}
当前 AI 安全领域普遍存在“AI 防御优势论”——
认为防御方部署大模型即可获得碾压式优势，而攻击方仍使用传统工具。
本文系统批判这一论调背后的三个隐含假设，
指出攻击方同样可以低成本启用 AI，且按需生成的攻击策略在成本结构上具有天然优势。
对抗样本攻防将博弈从网络层精确转移到模型参数空间，结构完全同构。
防御方引入 AI 后面临虚警消耗、模型退化和分析师能力丧失的三重陷阱。
本文经外部审查修正，补充了“按需生成”成本优势的边界条件、
分析师能力保留的形式化模型，以及攻防战场转移的精确坐标对应。
核心结论：AI 没有打破资源有限性约束，
防御方在 AI 时代的建设性策略是将 AI 的全部调度维度纳入贪心择优，
同时刻意保留人工研判以维持长期安全能力。
\end{abstract}

\tableofcontents

\section{引言}

大型语言模型和生成式 AI 正在快速渗透网络安全领域。
大量研究和产品宣传描绘了一幅图景：
防御方部署 AI 进行告警分析、威胁狩猎和自动响应，
攻击方仍在使用传统手段，防御方将获得碾压式优势。

这一叙事忽略了一个根本事实：攻击方同样可以启用 AI。
当攻防双方都武装了 AI，
资源有限性约束不会消失，博弈结构在数学上被完整保留，
只是战场的坐标发生了精确的平移。
本文对这一热潮进行系统性批判，
用已有的贪心择优理论框架分析 AI 化攻防的对称性，
并在批判基础上给出防御方在 AI 时代的建设性策略。

\section{“AI 防御优势论”的三个隐含假设}

当前 AI 安全乐观论调通常隐含以下假设：

\begin{enumerate}[label=(\arabic*)]
    \item \textbf{创新速度不对称}：防御方能更快训练和部署 AI，
          攻击方仍在使用传统工具。
    \item \textbf{数据壁垒}：防御方拥有完整的网络日志和告警数据，
          攻击方无法获得同等训练数据。
    \item \textbf{算力碾压}：防御方可以部署大规模推理集群，
          单次告警分析成本趋近于零。
\end{enumerate}

这些假设在攻击方不引入 AI 时成立。
一旦攻击方引入 AI，每一个假设都会被瓦解。

\section{攻击方的 AI 成本优势——按需生成对全量分析}

攻击方不一定需要自己训练模型。
开源大模型权重可被直接下载、量化后部署在普通 GPU 上，
暗网存在专门针对安全检测微调的模型交易。

攻击方用 AI 是“按需生成”，防御方是“全量分析”。
攻击方只需要生成一条能绕过检测的攻击 payload，
防御方需要用 AI 分析海量告警中的每一条。
两者的成本结构有本质差异。

设攻击方生成一条能逃避检测的 payload 的推理成本为 \(c_A^{\text{gen}}\)，
单次生成的成功概率为 \(p_{\text{evade}}\)，成功后的收益为 \(A_{\text{target}}\)。
攻击方仅在期望收益为正时使用生成模型。

\textbf{“按需生成”成本优势的边界条件}：
攻击方的有效成本不是 \(c_A^{\text{gen}}\)，而是 \(c_A^{\text{gen}} / p_{\text{evade}}\)——
攻击方可能需要生成多个变种才有一个成功。
若防御方通过持续对抗训练将 \(p_{\text{evade}}\) 压到极低，
攻击方的有效生成成本将逼近甚至超过防御方的单条告警分析成本。
此时“按需生成”的成本优势可能被逆转。
但防御方维持低 \(p_{\text{evade}}\) 的对抗训练成本本身也极高——
双方在 \(p_{\text{evade}}\) 和训练频率上形成新的均衡。

防御方用 AI 分析单条告警的推理成本为 \(c_D^{\text{analyze}}\)，
即使极低，也必须乘以海量告警基数 \(N_{\text{alerts}}\)，
总成本为 \(N_{\text{alerts}} \cdot c_D^{\text{analyze}}\)。
两者的有效成本结构有本质差异——
攻击方按需生成，防御方被海量告警基数锁死。
但这一差异的幅度依赖于 \(p_{\text{evade}}\) 的水平，
并非绝对的碾压关系。

\section{对抗样本军备竞赛的对称性}

对抗样本是 AI 化攻防中最纯粹的对称战场。
攻击方给恶意 payload 添加微小扰动使其逃逸 AI 检测，
防御方将这些对抗样本纳入训练集增强模型鲁棒性。

设攻击方生成一个对抗样本的成本为 \(c_A^{\text{adv}}\)，
该样本逃避防御方 AI 检测的概率为 \(p_{\text{evade}}^{(t)}\)。
防御方进行一轮对抗训练的成本为 \(c_D^{\text{adv\_train}}\)，
训练后 \(p_{\text{evade}}^{(t+1)}\) 下降 \(\Delta p_{\text{evade}}\)。

攻击方在 \(p_{\text{evade}}^{(t)} \cdot A_{\text{target}} > c_A^{\text{adv}}\) 时发起攻击。
防御方在 \(\Delta p_{\text{evade}} \cdot A_{\text{target}} \cdot N_{\text{adv\_alerts}} > c_D^{\text{adv\_train}}\) 时执行对抗训练。

双方在边际收益相等的点上达成均衡。

\section{战场转移的精确坐标}

AI 没有消除博弈，只是将战场从网络流量层精确转移到了模型参数空间。
以下是战场转移的精确坐标对应：

\begin{center}
\begin{tabular}{c|c}
\toprule
\textbf{传统战场坐标} & \textbf{AI 战场坐标} \\
\midrule
攻击方探测防御方分析率 & 攻击方查询防御方模型边界 \\
防御方周期性扰动分析参数 & 防御方周期性微调模型权重 \\
攻击方制造假告警逼移 & 攻击方制造对抗样本逃逸检测 \\
防御方反探测伪造响应 & 防御方模型输出叠加随机噪声 \\
攻击方联盟共享侦察信息 & 攻击方联盟共享模型查询和对抗样本 \\
\bottomrule
\end{tabular}
\end{center}

每一对战场坐标在数学结构上完全同构——
博弈的规则、成本和收益函数在新坐标系中被完整保留。

\section{数据投毒与模型窃取的柯克霍夫保护}

攻击方可能污染防御方训练数据来植入后门，
或通过大量查询来推断防御方模型边界。
模型权重和训练数据成为新的秘密参数，受柯克霍夫原则保护。

防御方模型架构可以公开，微调后的权重属于秘密层。
训练数据分布同样必须保密——
攻击方一旦掌握，便可针对性生成对抗样本。

周期性微调对应周期性扰动，
使攻击方通过历史查询推断的模型边界强制过期。
模型输出上叠加随机噪声，
使攻击方每次探测结果带有不可消除的不确定性。
这与 UEBA 基线参数的柯克霍夫保护完全同构。

\section{防御方的“AI 成本陷阱”}

防御方引入 AI 时面临三重相互关联的陷阱。

\textbf{AI 虚警消耗分析能力}。
AI 产生误报时，分析师的研判精力被无效告警消耗。
若 AI 误报率过高，净分析能力收益可能为负——
AI 节省的时间被自身误报吞噬殆尽。
这与 UEBA 误报负反馈循环完全同构。

\textbf{AI 模型退化}。
攻击方持续对抗训练导致防御方模型精度缓慢下降。
防御方被迫持续投入再训练成本，
形成“不训练就退化，训练就消耗算力”的两难困境。
最优再训练频率是边际精度维护收益等于边际训练成本的均衡点。

\textbf{AI 依赖削弱人的能力}。
分析师长期依赖 AI 裁定后，自身研判技能退化。
当 AI 模型因对抗攻击突然失效时，
分析师的研判能力不足以承接。

\textbf{人工研判保留比例的形式化}：
设人工研判保留比例为 \(r_{\text{manual}} \in [0,1]\)，
AI 辅助研判比例为 \(1 - r_{\text{manual}}\)。
\(r_{\text{manual}}\) 越高，分析师能力维持越好，但短期分析效率越低。
最优保留比例 \(r_{\text{manual}}^*\)
是边际长期能力维持收益等于边际短期效率损失的均衡点：
\[
\frac{d}{dr_{\text{manual}}} \left[ \text{AI失效时的承接能力} \right]
= \frac{d}{dr_{\text{manual}}} \left[ \text{当前分析效率损失} \right].
\]
该均衡条件在结构上与已有的所有贪心择优裁定同构。
\(r_{\text{manual}}^*\) 的具体数值依赖于 AI 失效概率的估计
和分析师技能退化速率，需要实战数据标定。

\section{AI 化攻防的核心结论与柯克霍夫原则的边界条件}

当攻防双方都部署 AI 时，优势不会自动归于防御方。
攻击方的 AI 按需生成在成本结构上具有天然优势，
但这一优势依赖于生成逃逸概率 \(p_{\text{evade}}\) 的水平，
双方在 \(p_{\text{evade}}\) 和训练频率上形成新的均衡。
AI 化攻防没有消除博弈，只是将战场从网络流量层
精确转移到了模型参数空间，结构在数学上完全保留。
防御方面临 AI 成本陷阱——虚警消耗、模型退化、分析师能力丧失，
三者叠加可能导致总安全能力不升反降。

\textbf{柯克霍夫原则在 AI 时代的边界条件}：
在第2卷和第3卷中，防御方通过柯克霍夫原则获得不对称优势——
参数保密和周期性扰动让攻击方的探测努力持续过期。
但在 AI 时代，如果攻击方可以通过大模型无限生成对抗样本，
攻击方的探测成本可能趋近于零。
周期性扰动对抗无限探测的有效性，
依赖于防御方的扰动频率是否能超过攻击方的探测生成速度。
当攻击方能够以趋近于零的成本持续生成探测时，
周期性扰动可能不足以消耗攻击方的探测预算。
这不是柯克霍夫原则的失效，
而是它的有效性边界需要在 AI 时代被重新审视。

\section{建设性框架：防御方在 AI 时代的策略}

在泼完冷水之后，本文给出防御方在 AI 时代的具体建设性策略。
防御方应将 AI 的模型选择、推理深度和对抗训练频率
全部纳入贪心择优的统一调度——
不是谁的 AI 更强，而是谁能在有限资源下把 AI 用得刚好够用。
同时，防御方必须刻意保留一定比例的人工研判，
维持分析师在 AI 失效时的承接能力。
这一保留比例本身由贪心择优裁定，
与其他分析维度的资源分配共享同一套经济逻辑。
最后，防御方应将模型权重、微调数据集和对抗训练策略
纳入柯克霍夫原则的秘密参数保护范围，
周期性微调和随机噪声共同对抗模型窃取和探测。

\section{诚实标注}

\begin{center}
\begin{xltabular}{\textwidth}{c|X|c}
\toprule
\textbf{命题} & \textbf{状态} & \textbf{层级} \\
\midrule
“AI 防御优势论”的三个隐含假设及其瓦解 &
逻辑推演，基于攻击方 AI 成本结构的分析 &
II（攻击方模型获取成本需实证数据） \\
攻击方的 AI 成本优势——按需生成对全量分析 &
推理成本比较逻辑成立，已补充 \(p_{\text{evade}}\) 修正 &
II（具体成本数值和逃逸概率需实测标定） \\
按需生成成本优势的边界条件 &
若防御方将 \(p_{\text{evade}}\) 压到极低，
攻击方有效成本将上升，双方在 \(p_{\text{evade}}\) 和训练频率上形成新均衡 &
II \\
对抗样本军备竞赛的对称性 &
攻防双方的成本收益模型已建立，与逼移博弈同构 &
II \\
战场转移的精确坐标 &
五对坐标映射已清晰列出，结构同构性已论证 &
II \\
数据投毒与模型窃取的柯克霍夫保护 &
与已有框架完全同构 &
II \\
防御方的 AI 成本陷阱——虚警、退化、能力丧失 &
三项陷阱的逻辑已建立 &
II（实战数据待验证） \\
人工研判保留比例的形式化 &
\(r_{\text{manual}}^*\) 的均衡条件已给出，与已有框架完全同构 &
II（具体数值需实测标定） \\
柯克霍夫原则在 AI 时代的边界条件 &
周期性扰动对抗无限探测的有效性依赖于扰动频率与探测生成速度之比 &
III（推测性分析，需进一步验证） \\
防御方在 AI 时代的建设性框架 &
批判基础上给出具体策略，逻辑自洽 &
II \\
\bottomrule
\end{xltabular}
\end{center}

\section{结语}

本文为当前 AI 安全狂热泼了一盆冷水。
AI 引入攻防双方后，资源有限性约束没有消失，
博弈结构被完整保留，只是战场的坐标发生了精确的平移。
攻击方按需生成的低成本优势、对抗样本军备竞赛的对称性、
防御方 AI 成本陷阱的三重困境——
共同指向一个结论：
AI 不会为任何一方带来碾压式优势。

防御方在 AI 时代的建设性策略是：
将 AI 的全部调度维度纳入贪心择优的统一框架，
同时刻意保留人工研判以防止人的能力退化。
够用就行原则在 AI 时代同样有效——
不是谁的 AI 更强，而是谁能在有限资源下把 AI 用得刚好够用。

\vspace{5mm}
\noindent\textbf{文档信息}：
\begin{itemize}
    \item 作者：李龙胜
    \item 版本：修正版（整合外部审查反馈）
    \item 许可：CC BY 4.0
    \item 前置依赖：四卷体系、AI 化攻防专题卷、UEBA 桥接、SOAR 剧本优化
\end{itemize}

\end{document}